\section{What Do the Agents' Scores Measure?}
\label{sec:scores}

A shortcut in a benchmark does not mean that the agents use it. ExCyTIn's authors publish the full logs of \lgModels{} agent runs on the 589 test questions behind their reported scores (ten models; GPT-5.1 at four reasoning settings and GPT-5 at two), with the judge's verdict for every question. These logs let us ask, question by question, \textbf{RQ3}: do the agents succeed because of the shortcut? We registered the analysis before reading the logs (\cref{sec:protocol}); success means that the judge accepted the final answer.

\begin{table*}[t]
  \caption{Agents do not succeed more often when the question names the answer's alert. Success (\%) of each published run on the 589 ExCyTIn test questions: all questions; questions that name the end alert (337) and the others (252), with the difference and its 95\% interval (incidents resampled); questions the alert-table lookup answers exactly (``Lookup'', 140); and the shortcut-resistant questions (``Resist.'', \lgResN: not named, and neither lookup answers them), with the run's rank on them. Runs are sorted by overall success.}
  \label{tab:logs}
  \centering\footnotesize
  % generated by paper/make_audit.py -- do not edit
\begin{tabular}{@{}lrrrrrr@{}}
\toprule
Model & All & Named & Other & $\Delta$ [95\% CI] & Lookup & Resist. \\
\midrule
Claude-Opus-4.5 & 59.8 & 61.7 & 57.1 & +4.6 \scriptsize[$-$4, +12] & 69.3 & 52.7 (2) \\
GPT-5.1 (high) & 57.6 & 57.3 & 57.9 & $-$0.7 \scriptsize[$-$9, +7] & 61.4 & 53.8 (1) \\
GPT-5 (high) & 55.3 & 55.8 & 54.8 & +1.0 \scriptsize[$-$9, +12] & 64.3 & 52.2 (4) \\
GPT-5.1 (medium) & 53.3 & 53.4 & 53.2 & +0.2 \scriptsize[$-$13, +13] & 50.7 & 52.7 (3) \\
Claude-Sonnet-4.5 & 47.7 & 50.4 & 44.0 & +6.4 \scriptsize[$-$3, +15] & 56.4 & 42.5 (5) \\
o3 & 44.7 & 46.0 & 42.9 & +3.1 \scriptsize[$-$9, +15] & 50.7 & 41.4 (6) \\
GPT-5.1 (low) & 44.3 & 46.6 & 41.3 & +5.3 \scriptsize[$-$7, +17] & 48.6 & 40.3 (7) \\
GPT-5 & 36.3 & 36.8 & 35.7 & +1.1 \scriptsize[$-$12, +14] & 40.7 & 33.9 (8) \\
GPT-5-mini & 35.8 & 35.6 & 36.1 & $-$0.5 \scriptsize[$-$8, +8] & 43.6 & 31.2 (10) \\
Grok-4 & 33.4 & 32.9 & 34.1 & $-$1.2 \scriptsize[$-$9, +5] & 34.3 & 31.2 (11) \\
GPT-5.1 (none) & 31.1 & 29.7 & 32.9 & $-$3.3 \scriptsize[$-$12, +5] & 32.1 & 32.8 (9) \\
Qwen3-235B-Thinking & 29.2 & 29.1 & 29.4 & $-$0.3 \scriptsize[$-$8, +6] & 32.1 & 29.0 (12) \\
Claude-Haiku-4.5 & 19.4 & 19.3 & 19.4 & $-$0.2 \scriptsize[$-$8, +8] & 22.9 & 18.3 (13) \\
GPT-5-nano & 8.7 & 8.9 & 8.3 & +0.6 \scriptsize[$-$4, +6] & 11.4 & 8.6 (14) \\
\bottomrule
\end{tabular}

\end{table*}

\paragraph{Naming the alert does not help the agents} If agents found the answer by finding the named alert, they would succeed more often on named questions. They do not (\cref{tab:logs}). Only \lgPos{} of the \lgModels{} runs succeed more often on named questions (one-sided sign test $p=\lgP$); the differences range from \lgDmin{} to \lgDmax{} points, and no interval excludes zero. The per-incident correlation of \cref{sec:excytin} therefore comes from differences between incidents, not from the naming. Questions that the alert-table lookup answers exactly are somewhat easier: \lgTPos{} of \lgModels{} runs do better on them, by \lgTmin{} to \lgTmax{} points (median \lgTmedian), but only \lgTCI{} of these differences have intervals that exclude zero, and these questions make up a share of each run's successes (\lgShareMin--\lgShareMax\%) close to their share of the test set (\oTable\%).

\paragraph{The agents investigate even when one query would do} The logs show why the shortcut is not used. On the 140 questions that a single query of the alert table answers, Claude-Opus-4.5 issues \tcOpusStepsLookup{} SQL queries on average, against \tcOpusStepsOther{} on the other questions; GPT-5 issues \tcGFiveStepsLookup{} and \tcGFiveStepsOther, o3 \tcOThreeStepsLookup{} and \tcOThreeStepsOther. On these questions, no run among the three queries the \texttt{SecurityAlert} table in more than \tcMaxAlertTable\% of cases (Claude-Opus-4.5: \tcOpusAlertTable\%). The agents follow the generic strategy of the benchmark's prompt, exploring the schema and pivoting through event tables, whatever the question says.

\paragraph{Shortcut-resistant questions rank the agents the same way} On the \lgResN{} questions that neither name the end alert nor yield to either lookup, every run scores at most \lgDropMax{} points below its overall success, and the ranking of the runs is nearly unchanged (Kendall $\tau=\lgTau$; \cref{fig:logs}). The four strongest runs swap neighboring places, and no run moves by more than two ranks.

\begin{figure}[t]
  \centering
  \includegraphics[width=\columnwidth]{fig_logs.pdf}
  \caption{Success depends only weakly on how much investigation a question needs. For each published run: success on all questions, on the 140 questions one alert-table lookup answers, and on the \lgResN{} shortcut-resistant questions.}
  \label{fig:logs}
\end{figure}

\paragraph{Depth explains little of the failures} The best run fails \tcOpusFail{} of the 140 questions that one lookup answers; it submits a wrong answer in all but \tcOpusFailNoAns{} of them, so it neither runs out of steps nor gives up. In an exploratory split by the length of the generation path, which ExCyTIn uses as its measure of difficulty, Claude-Opus-4.5 succeeds on \plOpusOne\% of the \plN{} questions whose answer lies in the alert already described in the context, against \plOpusThree\% for paths of length three and \plOpusFive\% for longer ones. Even with nothing to investigate, the strongest agent misses more than a quarter of the answers.

\smallskip\noindent\textbf{Finding 2 (RQ3).} The ExCyTIn agents do not exploit the naming shortcut: their success is the same whether or not the question names the answer's alert, and the ranking of strong agents survives on shortcut-resistant questions. What the score reflects only weakly is the depth of investigation: a question that one lookup answers is a few points easier than one that needs a multi-hop pivot, and the strongest agent misses more than a quarter of the questions with nothing to investigate. The benchmark measures how reliably an agent retrieves and reports the right value more than how deeply it investigates.
